\chapter{Communication Complexity — TV Distance}
%%% generated by the blueprint pipeline from TCSlib.CommunicationComplexity.NewmanTheorem.TVDistance
\section{Overview}

This module develops the theory of total variation (TV) distance between probability
measures on measurable spaces.  It gives two equivalent definitions---one via the total
variation norm of a signed measure and one as a supremum over measurable events---proves
their equivalence, establishes the basic bound on event probability gaps, and derives a
concrete half-$\ell^1$ formula on finite spaces, specialised further to Boolean and
product distributions.

%%%%%%%%%%%%%%%%%%%%%%%%%%%%%%%%%%%%%%%%%%%%%%%%%%%%%%%%%%%%%%%%%%%%%%%%%%%%%%%
\section{Declarations}

\begin{definition}[Signed measure difference]
\lean{CommunicationComplexity.signedMeasureDiff}
\label{CommunicationComplexity.signedMeasureDiff}
\leanok
Given two probability measures $\mu$ and $\nu$ on a measurable space $\Omega$, the
\emph{signed measure difference} $\mu - \nu$ is the signed measure on $\Omega$ obtained
by taking the difference of $\mu$ and $\nu$ as measures.
\end{definition}

\begin{definition}[Total variation distance]
\lean{CommunicationComplexity.tvDistance}
\label{CommunicationComplexity.tvDistance}
\leanok
\uses{CommunicationComplexity.signedMeasureDiff}
The \emph{total variation distance} between two probability measures $\mu$ and $\nu$ on
$\Omega$ is defined as $\tfrac{1}{2}$ times the total variation norm of the signed
measure $\mu - \nu$, i.e.\
$\mathrm{TV}(\mu,\nu) = \tfrac{1}{2}\|\mu - \nu\|_{\mathrm{TV}}$.
\end{definition}

\begin{definition}[Total variation distance via event supremum]
\lean{CommunicationComplexity.tvDistanceSup}
\label{CommunicationComplexity.tvDistanceSup}
\leanok
The \emph{supremum definition} of total variation distance between $\mu$ and $\nu$ is
\[
  \mathrm{TV}_{\sup}(\mu,\nu)
  \;=\; \sup\bigl\{|\mu(S) - \nu(S)| : S \subseteq \Omega \text{ measurable}\bigr\}.
\]
\end{definition}

\begin{lemma}[Signed measure via Jordan decomposition]
\lean{CommunicationComplexity.TVDistance.signedMeasure_apply_eq_posPart_sub_negPart}
\label{CommunicationComplexity.TVDistance.signedMeasure_apply_eq_posPart_sub_negPart}
\leanok
\difficulty{2}
Let $s$ be a signed measure on a measurable space $\Omega$, and let $s = s^+ - s^-$ be
its Jordan decomposition into a pair of (finite, mutually singular) positive measures
$s^+$ and $s^-$. Then for every measurable set $S \subseteq \Omega$,
\[
s(S) = s^+(S) - s^-(S),
\]
where $s^+(S)$ and $s^-(S)$ denote the real-valued measures of $S$ under the positive
and negative parts.
\end{lemma}

\begin{lemma}[Signed measure difference vanishes on the whole space]
\lean{CommunicationComplexity.TVDistance.signedMeasureDiff_univ}
\label{CommunicationComplexity.TVDistance.signedMeasureDiff_univ}
\leanok
\uses{CommunicationComplexity.signedMeasureDiff}
\difficulty{1}
Let $\mu$ and $\nu$ be probability measures on a measurable space $\Omega$, and let $\mu
- \nu$ denote their signed measure difference. Then $(\mu - \nu)(\Omega) = 0$.
\end{lemma}

\begin{lemma}[Jordan parts of a probability-measure difference have equal total mass]
\lean{CommunicationComplexity.TVDistance.jordan_posPart_real_univ_eq_negPart_real_univ}
\label{CommunicationComplexity.TVDistance.jordan_posPart_real_univ_eq_negPart_real_univ}
\leanok
\uses{CommunicationComplexity.TVDistance.signedMeasureDiff_univ, CommunicationComplexity.TVDistance.signedMeasure_apply_eq_posPart_sub_negPart, CommunicationComplexity.signedMeasureDiff}
\difficulty{2}
Let $\mu$ and $\nu$ be probability measures on a measurable space $\Omega$, and let $\mu
- \nu$ be their signed measure difference. Writing $(\mu-\nu)^{+}$ and $(\mu-\nu)^{-}$
for the positive and negative parts of the Jordan decomposition of $\mu - \nu$, the two
parts assign the same total real mass to the whole space:
\[
(\mu-\nu)^{+}(\Omega) = (\mu-\nu)^{-}(\Omega).
\]
\end{lemma}

\begin{lemma}[Total variation mass as the sum of Jordan parts]
\lean{CommunicationComplexity.TVDistance.totalVariation_real_univ}
\label{CommunicationComplexity.TVDistance.totalVariation_real_univ}
\leanok
\difficulty{1}
Let $s$ be a signed measure on a measurable space $\Omega$, with Jordan decomposition $s
= s^+ - s^-$ into its positive and negative parts $s^+$ and $s^-$. Then the total mass
of the total variation measure $\abs{s}$ of $s$, taken as a real number, is the sum of
the total masses of the two Jordan parts:
\[
\abs{s}(\Omega) = s^+(\Omega) + s^-(\Omega).
\]
\end{lemma}

\begin{lemma}[Half the total variation equals the positive Jordan part]
\lean{CommunicationComplexity.TVDistance.half_totalVariation_real_univ_eq_posPart_real_univ}
\label{CommunicationComplexity.TVDistance.half_totalVariation_real_univ_eq_posPart_real_univ}
\leanok
\uses{CommunicationComplexity.TVDistance.jordan_posPart_real_univ_eq_negPart_real_univ, CommunicationComplexity.TVDistance.totalVariation_real_univ, CommunicationComplexity.signedMeasureDiff}
\difficulty{2}
Let $\mu$ and $\nu$ be probability measures on a measurable space $\Omega$, and let
$\mu-\nu$ denote their signed measure difference. Writing $\abs{\mu-\nu}$ for its total
variation measure and $(\mu-\nu)^{+}$ for the positive part of its Jordan decomposition,
one has
\[
\tfrac{1}{2}\,\abs{\mu-\nu}(\Omega) = (\mu-\nu)^{+}(\Omega),
\]
where both sides are evaluated as real numbers on the whole space $\Omega$.
\end{lemma}

\begin{lemma}[Event probability gap bounded by half the total variation]
\lean{CommunicationComplexity.TVDistance.event_abs_signedMeasureDiff_le_half_totalVariation}
\label{CommunicationComplexity.TVDistance.event_abs_signedMeasureDiff_le_half_totalVariation}
\leanok
\uses{CommunicationComplexity.TVDistance.half_totalVariation_real_univ_eq_posPart_real_univ, CommunicationComplexity.TVDistance.jordan_posPart_real_univ_eq_negPart_real_univ, CommunicationComplexity.TVDistance.signedMeasure_apply_eq_posPart_sub_negPart, CommunicationComplexity.signedMeasureDiff}
\difficulty{4}
Let $\mu$ and $\nu$ be probability measures on a measurable space $\Omega$, and let $\mu
- \nu$ denote their signed measure difference. Then for every measurable set $S
\subseteq \Omega$,
\[
  \abs{(\mu - \nu)(S)} \;\le\; \tfrac{1}{2}\, \abs{\mu - \nu}(\Omega),
\]
where $\abs{\mu - \nu}$ is the total variation measure of $\mu - \nu$ and $\abs{\mu -
\nu}(\Omega)$ is its total mass.
\end{lemma}

\begin{lemma}[Probability gap bounded by half total variation]
\lean{CommunicationComplexity.TVDistance.event_abs_measureReal_sub_le_half_totalVariation}
\label{CommunicationComplexity.TVDistance.event_abs_measureReal_sub_le_half_totalVariation}
\leanok
\uses{CommunicationComplexity.TVDistance.event_abs_signedMeasureDiff_le_half_totalVariation, CommunicationComplexity.signedMeasureDiff}
\difficulty{2}
Let $\mu$ and $\nu$ be probability measures on a measurable space $\Omega$, and let $\mu
- \nu$ denote their signed measure difference. Then for every measurable set $S
\subseteq \Omega$,
\[
\abs{\mu(S) - \nu(S)} \le \tfrac{1}{2}\,\abs{\mu - \nu}(\Omega),
\]
where $\abs{\mu - \nu}$ is the total variation measure of $\mu - \nu$.
\end{lemma}

\begin{lemma}[Attainment of half the total variation by a single event]
\lean{CommunicationComplexity.TVDistance.exists_event_abs_signedMeasureDiff_eq_half_totalVariation}
\label{CommunicationComplexity.TVDistance.exists_event_abs_signedMeasureDiff_eq_half_totalVariation}
\leanok
\uses{CommunicationComplexity.TVDistance.half_totalVariation_real_univ_eq_posPart_real_univ, CommunicationComplexity.TVDistance.signedMeasure_apply_eq_posPart_sub_negPart, CommunicationComplexity.signedMeasureDiff}
\difficulty{4}
Let $\mu$ and $\nu$ be probability measures on a measurable space $\Omega$, and let $\mu
- \nu$ denote their signed measure difference. Then there exists a measurable set $S
\subseteq \Omega$ at which
\[
\abs{(\mu - \nu)(S)} = \tfrac{1}{2}\,\abs{\mu - \nu}(\Omega),
\]
where $\abs{\mu - \nu}$ is the total variation measure of $\mu - \nu$; that is, some
event realizes exactly half the total variation of $\mu - \nu$ over the whole space.
\end{lemma}

\begin{lemma}[Attainment of the total-variation distance]
\lean{CommunicationComplexity.TVDistance.exists_event_abs_measureReal_sub_eq_half_totalVariation}
\label{CommunicationComplexity.TVDistance.exists_event_abs_measureReal_sub_eq_half_totalVariation}
\leanok
\uses{CommunicationComplexity.TVDistance.exists_event_abs_signedMeasureDiff_eq_half_totalVariation, CommunicationComplexity.signedMeasureDiff}
\difficulty{2}
Let $\mu$ and $\nu$ be probability measures on a measurable space $\Omega$, and let $\mu
- \nu$ denote their signed measure difference, with total variation measure
$\abs{\mu-\nu}$. Then there exists a measurable set $S \subseteq \Omega$ such that
\[
  \abs{\mu(S) - \nu(S)} = \tfrac{1}{2}\,\abs{\mu-\nu}(\Omega).
\]
\end{lemma}

\begin{theorem}[Equivalence of the two definitions of total variation distance]
\lean{CommunicationComplexity.TVDistance.tvDistance_eq_tvDistanceSup}
\label{CommunicationComplexity.TVDistance.tvDistance_eq_tvDistanceSup}
\leanok
\uses{CommunicationComplexity.TVDistance.event_abs_measureReal_sub_le_half_totalVariation, CommunicationComplexity.TVDistance.exists_event_abs_measureReal_sub_eq_half_totalVariation, CommunicationComplexity.signedMeasureDiff, CommunicationComplexity.tvDistance, CommunicationComplexity.tvDistanceSup}
\difficulty{4}
Let $\mu$ and $\nu$ be probability measures on a measurable space $\Omega$. Then the
total variation distance $\mathrm{TV}(\mu,\nu) =
\tfrac{1}{2}\norm{\mu-\nu}_{\mathrm{TV}}$, defined as half the total variation norm of
the signed measure $\mu-\nu$, agrees with its supremum form over events,
\[
\mathrm{TV}(\mu,\nu) \;=\; \sup\bigl\{\,\abs{\mu(S)-\nu(S)} : S \subseteq \Omega \text{
measurable}\,\bigr\}.
\]
\end{theorem}

\begin{theorem}[Events differ by at most the total variation distance]
\lean{CommunicationComplexity.TVDistance.abs_measureReal_sub_le_tvDistance}
\label{CommunicationComplexity.TVDistance.abs_measureReal_sub_le_tvDistance}
\leanok
\uses{CommunicationComplexity.TVDistance.event_abs_measureReal_sub_le_half_totalVariation, CommunicationComplexity.tvDistance}
\difficulty{1}
Let $\mu$ and $\nu$ be probability measures on a measurable space $\Omega$. Then for
every measurable set $S \subseteq \Omega$,
\[
  \abs{\mu(S) - \nu(S)} \le \mathrm{TV}(\mu,\nu),
\]
where $\mathrm{TV}(\mu,\nu) = \tfrac{1}{2}\norm{\mu - \nu}_{\mathrm{TV}}$ is the total
variation distance between $\mu$ and $\nu$.
\end{theorem}

\begin{theorem}[Nonnegativity of total variation distance]
\lean{CommunicationComplexity.TVDistance.tvDistance_nonneg}
\label{CommunicationComplexity.TVDistance.tvDistance_nonneg}
\leanok
\uses{CommunicationComplexity.tvDistance}
\difficulty{1}
Let $\mu$ and $\nu$ be probability measures on a measurable space $\Omega$. Then their
total variation distance is nonnegative, $\mathrm{TV}(\mu,\nu) \ge 0$.
\end{theorem}

\begin{lemma}[Indicator sum bounded by the positive-part sum]
\lean{CommunicationComplexity.TVDistance.sum_indicator_le_sum_posPart}
\label{CommunicationComplexity.TVDistance.sum_indicator_le_sum_posPart}
\leanok
\difficulty{3}
Let $\alpha$ be a finite type and let $a : \alpha \to \bbr$ be a real-valued function.
For every subset $S \subseteq \alpha$,
\[
\sum_{x \in S} a(x) \;\le\; \sum_{x \in \alpha} \max\bigl(a(x),\,0\bigr).
\]
\end{lemma}

\begin{lemma}[Positive part sum equals half the $\ell^1$ norm]
\lean{CommunicationComplexity.TVDistance.sum_posPart_eq_half_sum_abs}
\label{CommunicationComplexity.TVDistance.sum_posPart_eq_half_sum_abs}
\leanok
\difficulty{4}
Let $a\colon\alpha\to\bbr$ be a real-valued function on a finite type $\alpha$ whose
values sum to zero, that is, $\sum_{x\in\alpha} a(x) = 0$. Then the sum of the positive
parts of $a$ equals half the sum of its absolute values:
\[
\sum_{x\in\alpha} \max\bigl(a(x),\,0\bigr) \;=\; \frac{1}{2}\sum_{x\in\alpha}
\abs{a(x)}.
\]
\end{lemma}

\begin{lemma}[Half-total-variation bound for a zero-sum indicator sum]
\lean{CommunicationComplexity.TVDistance.abs_sum_indicator_le_half_sum_abs}
\label{CommunicationComplexity.TVDistance.abs_sum_indicator_le_half_sum_abs}
\leanok
\uses{CommunicationComplexity.TVDistance.sum_indicator_le_sum_posPart, CommunicationComplexity.TVDistance.sum_posPart_eq_half_sum_abs}
\difficulty{4}
Let $\alpha$ be a finite type and let $a : \alpha \to \bbr$ be a real-valued function
whose entries sum to zero, $\sum_{x \in \alpha} a(x) = 0$. Then for every subset $S
\subseteq \alpha$,
\[
  \Bigl|\sum_{x \in S} a(x)\Bigr| \;\le\; \tfrac{1}{2} \sum_{x \in \alpha} \abs{a(x)}.
\]
\end{lemma}

\begin{lemma}[Positive part of a zero-sum family equals half its total variation]
\lean{CommunicationComplexity.TVDistance.abs_sum_pos_indicator_eq_half_sum_abs}
\label{CommunicationComplexity.TVDistance.abs_sum_pos_indicator_eq_half_sum_abs}
\leanok
\uses{CommunicationComplexity.TVDistance.sum_posPart_eq_half_sum_abs}
\difficulty{2}
Let $\alpha$ be a finite type and let $a : \alpha \to \bbr$ be a family of real numbers
whose total sum vanishes, $\sum_{x} a(x) = 0$. Then
\[
\Bigl|\sum_{x} \max(a(x), 0)\Bigr| = \tfrac{1}{2} \sum_{x} \abs{a(x)}.
\]
\end{lemma}

\begin{definition}[Singleton mass difference]
\lean{CommunicationComplexity.TVDistance.singletonMassDiff}
\label{CommunicationComplexity.TVDistance.singletonMassDiff}
\leanok
For probability measures $\mu$ and $\nu$ on $\Omega$ and a point $\omega \in \Omega$,
the \emph{singleton mass difference} is $\mu(\{\omega\}) - \nu(\{\omega\}) \in \mathbb{R}$.
\end{definition}

\begin{lemma}[Probability gap as a sum of singleton mass differences]
\lean{CommunicationComplexity.TVDistance.measureReal_sub_eq_sum_indicator_singletonMassDiff}
\label{CommunicationComplexity.TVDistance.measureReal_sub_eq_sum_indicator_singletonMassDiff}
\leanok
\uses{CommunicationComplexity.FiniteMeasureSpace, CommunicationComplexity.FiniteMeasureSpace.measureReal_eq_sum_singletons, CommunicationComplexity.TVDistance.singletonMassDiff}
\difficulty{3}
Let $\Omega$ be a finite discrete measurable space, and let $\mu$ and $\nu$ be
probability measures on $\Omega$. Then for every subset $S \subseteq \Omega$,
\[
\mu(S) - \nu(S) = \sum_{\omega \in \Omega} \mathbf{1}[\omega \in
S]\,\bigl(\mu(\{\omega\}) - \nu(\{\omega\})\bigr),
\]
where the sum ranges over all points of $\Omega$ and each term contributes the singleton
mass difference $\mu(\{\omega\}) - \nu(\{\omega\})$ exactly when $\omega \in S$ and $0$
otherwise.
\end{lemma}

\begin{lemma}[Singleton mass differences sum to zero]
\lean{CommunicationComplexity.TVDistance.sum_singletonMassDiff_eq_zero}
\label{CommunicationComplexity.TVDistance.sum_singletonMassDiff_eq_zero}
\leanok
\uses{CommunicationComplexity.FiniteMeasureSpace, CommunicationComplexity.FiniteMeasureSpace.measureReal_eq_sum_singletons, CommunicationComplexity.TVDistance.singletonMassDiff}
\difficulty{3}
Let $\Omega$ be a finite discrete measurable space, and let $\mu$ and $\nu$ be
probability measures on $\Omega$. Then the singleton mass differences sum to zero:
\[
\sum_{\omega \in \Omega} \bigl(\mu(\{\omega\}) - \nu(\{\omega\})\bigr) = 0.
\]
\end{lemma}

\begin{theorem}[Total variation distance as half the $\ell^1$ distance]
\lean{CommunicationComplexity.TVDistance.tvDistanceSup_eq_half_sum}
\label{CommunicationComplexity.TVDistance.tvDistanceSup_eq_half_sum}
\leanok
\uses{CommunicationComplexity.FiniteMeasureSpace, CommunicationComplexity.TVDistance.abs_sum_indicator_le_half_sum_abs, CommunicationComplexity.TVDistance.abs_sum_pos_indicator_eq_half_sum_abs, CommunicationComplexity.TVDistance.measureReal_sub_eq_sum_indicator_singletonMassDiff, CommunicationComplexity.TVDistance.singletonMassDiff, CommunicationComplexity.TVDistance.sum_singletonMassDiff_eq_zero, CommunicationComplexity.tvDistanceSup}
\difficulty{5}
Let $\Omega$ be a finite set equipped with the discrete measurable structure, and let
$\mu$ and $\nu$ be probability measures on $\Omega$. Then the supremum form of total
variation distance, $\mathrm{TV}_{\sup}(\mu,\nu) = \sup\{|\mu(S) - \nu(S)| : S \subseteq
\Omega\}$, is exactly half the $\ell^1$ distance between the two measures across the
individual points:
\[
\mathrm{TV}_{\sup}(\mu,\nu) \;=\; \frac{1}{2}\sum_{\omega \in \Omega}
\bigl|\mu(\{\omega\}) - \nu(\{\omega\})\bigr|.
\]
\end{theorem}

\begin{theorem}[Total variation distance as half the $\ell^1$ distance]
\lean{CommunicationComplexity.TVDistance.tvDistance_eq_half_sum}
\label{CommunicationComplexity.TVDistance.tvDistance_eq_half_sum}
\leanok
\uses{CommunicationComplexity.FiniteMeasureSpace, CommunicationComplexity.TVDistance.tvDistanceSup_eq_half_sum, CommunicationComplexity.TVDistance.tvDistance_eq_tvDistanceSup, CommunicationComplexity.tvDistance}
\difficulty{1}
Let $\Omega$ be a finite discrete measurable space, and let $\mu$ and $\nu$ be
probability measures on $\Omega$. Then the total variation distance between $\mu$ and
$\nu$ equals half the sum, over all points $\omega \in \Omega$, of the absolute
differences of their point masses:
\[
  \mathrm{TV}(\mu,\nu) \;=\; \frac{1}{2}\sum_{\omega \in \Omega}
    \abs{\mu(\{\omega\}) - \nu(\{\omega\})}.
\]
\end{theorem}

\begin{theorem}[Total variation distance on the two-point space]
\lean{CommunicationComplexity.TVDistance.tvDistance_bool_eq_abs_true}
\label{CommunicationComplexity.TVDistance.tvDistance_bool_eq_abs_true}
\leanok
\uses{CommunicationComplexity.TVDistance.tvDistance_eq_half_sum, CommunicationComplexity.tvDistance}
\difficulty{4}
Let $\mu$ and $\nu$ be probability measures on the two-point space $\mathrm{Bool} =
\{\mathrm{true}, \mathrm{false}\}$. Then their total variation distance is the absolute
difference of the masses they assign to the single point $\mathrm{true}$:
\[
\mathrm{TV}(\mu,\nu) = \abs{\mu(\{\mathrm{true}\}) - \nu(\{\mathrm{true}\})}.
\]
\end{theorem}

\begin{theorem}[Subadditivity of total variation distance under products]
\lean{CommunicationComplexity.TVDistance.tvDistance_prod_le}
\label{CommunicationComplexity.TVDistance.tvDistance_prod_le}
\leanok
\uses{CommunicationComplexity.FiniteMeasureSpace, CommunicationComplexity.TVDistance.tvDistance_eq_half_sum, CommunicationComplexity.tvDistance}
\difficulty{5}
Let $\alpha$ and $\beta$ be finite discrete measurable spaces, let $\mu_1,\nu_1$ be
probability measures on $\alpha$, and let $\mu_2,\nu_2$ be probability measures on
$\beta$. Writing $\mathrm{TV}$ for the total variation distance and $\mu_1\otimes\mu_2$
for the product measure, the total variation distance between the two product measures
satisfies
\[
\mathrm{TV}(\mu_1\otimes\mu_2,\ \nu_1\otimes\nu_2)\ \le\
\mathrm{TV}(\mu_1,\nu_1)+\mathrm{TV}(\mu_2,\nu_2).
\]
\end{theorem}
